\documentclass[11pt]{article}
\usepackage[T1]{fontenc}
\usepackage{lmodern}
\usepackage[margin=1.05in]{geometry}
\usepackage{amsmath,amssymb,amsthm,mathtools}
\usepackage{tikz-cd}
\usepackage{microtype}
\usepackage[colorlinks=true,linkcolor=blue,citecolor=blue,urlcolor=blue]{hyperref}
\hypersetup{pdftitle={Maximal topological complexity of connected sums of aspherical manifolds},pdfauthor={Kaiwen Jin}}
\newtheorem{theorem}{Theorem}[section]
\newtheorem{proposition}[theorem]{Proposition}
\newtheorem{lemma}[theorem]{Lemma}
\theoremstyle{remark}
\newtheorem{remark}[theorem]{Remark}
\newcommand{\TC}{\operatorname{TC}}
\newcommand{\Q}{\mathbb{Q}}
\newcommand{\Z}{\mathbb{Z}}
\newcommand{\im}{\operatorname{im}}
\newcommand{\bigsq}{\bigsqcup}
\title{Maximal topological complexity of connected sums\\of aspherical manifolds}
\author{Kaiwen Jin\\\small\href{mailto:yui.ui.siki@gmail.com}{yui.ui.siki@gmail.com}}
\date{October 13, 2026}
\begin{document}
\maketitle

\begin{abstract}
Let \(M\) and \(N\) be closed, connected, oriented aspherical manifolds of the
same dimension \(n\geq 3\). We prove that their connected sum has ordinary
unreduced topological complexity \(2n+1\). Thus joining two aspherical
summands always attains the dimensional upper bound, although the connected
sum is not required to be aspherical. The proof detects the two ordered
cross products of the summands' fundamental classes in a quotient of the
left-right join of their free-product fundamental group. A decomposition
into factor and cyclic cosets shows that these mixed classes survive, and
their survival obstructs a motion planner with \(2n\) domains.
\end{abstract}

\section{Introduction}

The topological complexity of a path-connected space measures how many
continuous rules are needed to choose a path between arbitrary endpoints.
Following Farber~\cite{Farber}, we use the unreduced convention: \(\TC(X)\)
is the least number of open sets covering \(X\times X\) on which
\[
 e:X^{[0,1]}\longrightarrow X\times X,\qquad
 e(\gamma)=(\gamma(0),\gamma(1)),
\]
has continuous sections. The path space has the compact-open topology.
For a connected \(n\)-manifold the dimensional upper bound is \(2n+1\).

Connected sums offer a natural test of how motion-planning complexity
changes under gluing. For aspherical summands the fundamental group of the
sum is a free product, and the corresponding classifying space is a wedge.
The free-product formula of Dranishnikov and Sadykov~\cite{DSfree} gives
\(\TC(M\vee N)=2n+1\) when \(M,N\) are closed oriented aspherical
\(n\)-manifolds, \(n\geq3\). This computation does not give a lower bound
for the connected sum: the fundamental-group isomorphism induced by
pinching the neck does not make the connected sum aspherical.

Neofytidis~\cite[Question 5.5(b)]{Neofytidis} asks whether
\(\TC(M\#N)=2n+1\) always holds in this setting. His Corollary 1.3
establishes the maximum for certain four-dimensional sums with a negatively
curved summand. The revised work of Daundkar, Santhanam and
Thandar~\cite[Proposition 5.3]{DST} computes the value \(7\) for wedges
of closed oriented aspherical three-manifolds and gives bounds for
specified connected sums. The question for connected sums therefore
requires an obstruction that survives on their own fundamental class.

\begin{theorem}\label{thm:main}
For every pair of closed, connected, oriented aspherical manifolds \(M,N\)
of the same dimension \(n\geq3\),
\[
 \TC(M\#N)=2n+1.
\]
\end{theorem}

The theorem includes summands with nontrivial center and aspherical
rational homology spheres. There is no curvature or intermediate
cohomology hypothesis. In particular, it gives
\(\TC(M\#M)=7\) for every closed oriented aspherical three-manifold \(M\),
including the torus and the examples in~\cite[Question 5.10]{Neofytidis}.

The argument uses the diagonal covering and left-right join construction
of Farber, Grant, Lupton and Oprea~\cite{FGLO}. We study its orbit space
rather than transfer the wedge computation directly. In the free product,
every simplex with nontrivial stabilizer has all its group labels in a
single factor coset or a hyperbolic cyclic coset. Distinct such cosets
intersect in at most one label. Their Borel constructions can contribute
only the two same-summand products in top homology. The two ordered mixed
products consequently survive in the orbit space, where a planner with
\(2n\) domains would force them to vanish.

We first describe this homological obstruction, then prove the coset
decomposition and its homology consequence. The final section applies it
to the connected-sum fundamental class. All homology below has coefficients
\(\Q\). Products and infinite CW constructions use the compactly generated
convention; see~\cite[Chapters 5 and 10]{May}.

\section{A homological obstruction from the diagonal covering}

Let \(G\) be a discrete group and put \(P=G\times G\). For a nonempty set
\(S\), let
\[
 E(S)=S*S*S*\cdots
\]
be the join of countably many copies of \(S\). Its vertices are
\((i,s)\), \(i\geq1\), and a simplex contains at most one vertex at each
level \(i\). We give it the weak CW topology. This infinite join is
contractible: a sphere map has compact image in a finite subcomplex, and
that image can be coned to a vertex at an unused level. The CW Whitehead
theorem then applies.

Set \(E=E(G)\). The left-right action
\[
 (u,v)\cdot(i,g)=(i,ugv^{-1})
\]
makes \(E\) a \(P\)-CW complex. An element that stabilizes a simplex fixes
it pointwise because it preserves the levels. Write \(E_m\) for the join
of the first \(m\) copies of \(G\), of dimension \(m-1\).
These are the join models in~\cite[Sections 2.5--2.9]{FGLO}.

Choose a contractible free \(P\)-CW complex \(EP\). For a \(P\)-space \(W\)
write \(W_{hP}=(EP\times W)/P\), with diagonal action, for its Borel
construction. The maps
\[
 \pi:E_{hP}\longrightarrow BP=EP/P,\qquad
 \rho:E_{hP}\longrightarrow E/P
\]
forget the \(E\) and \(EP\) coordinates, respectively. The first is an
associated bundle with contractible fiber, so it induces a homology
isomorphism. We shall use
\begin{equation}\label{eq:tau}
 \tau_d=\rho_*(\pi_*)^{-1}:
 H_d(P)\longrightarrow H_d(E/P),
\end{equation}
identifying group homology with the homology of \(BP\).

\begin{lemma}\label{lem:planner}
Let \(Y\) be a closed connected manifold with fundamental group \(G\),
and let \(c_Y:Y\to BG\) be a classifying map. If \(\TC(Y)\leq m\), then
for every \(d\geq m\) and \(z\in H_d(Y\times Y)\),
\[
 \tau_d\big((c_Y\times c_Y)_*z\big)=0.
\]
Here \(H_d(BG\times BG)\) is identified with \(H_d(P)\) through the
classifying homotopy equivalence inducing the identity on \(P\).
\end{lemma}

\begin{proof}
Let \(\widetilde Y\to Y\) be the universal cover. Lifting a path and taking
its two lifted endpoints modulo simultaneous deck transformations gives
a continuous map from \(Y^{[0,1]}\) to the diagonal covering
\[
 q:(\widetilde Y\times\widetilde Y)/\Delta G
       \longrightarrow Y\times Y,
\]
over the endpoint map. The map is independent of the chosen lift.
Continuity follows from local trivializations and path lifting; the
compact interval can be covered by finitely many trivializing
neighborhoods along each path. Thus every local motion planner gives a
local section of \(q\).

The covering \(q\) is associated to the principal \(P\)-cover
\(\widetilde Y\times\widetilde Y\to Y\times Y\) with fiber \(G\)
and action \(ugv^{-1}\); compare~\cite[proof of Theorem 2.1]{FGLO}.
Choose a partition of unity for the \(m\) section domains, with each
closed support contained in its domain. The weighted join of the local
sections is a section of the associated bundle with fiber \(E_m\).
It is continuous in the weak join topology: near a point, each nonzero
section label is locally constant in a covering trivialization, while
the weights whose supports miss that point vanish on a neighborhood.
The section is equivalently a \(P\)-equivariant map
\[
 f:\widetilde Y\times\widetilde Y\longrightarrow E_m.
\]

Choose a \(P\)-equivariant classifying map
\(j:\widetilde Y\times\widetilde Y\to EP\).
The map
\[
 b:Y\times Y\longrightarrow E_{hP},\qquad
 b([y_0,y_1])=[j(y_0,y_1),f(y_0,y_1)]
\]
projects to the classifying map of the principal \(P\)-cover, homotopic
to \(c_Y\times c_Y\) after the stated identification. Hence
\(\pi_*b_*=(c_Y\times c_Y)_*\). Its composition with \(\rho\) factors
through \(E_m/P\), a CW complex of dimension \(m-1\).
It vanishes on \(H_d\) for \(d\geq m\), which proves the lemma by
\eqref{eq:tau}. The universal cover of \(Y\) was not assumed contractible.
\end{proof}

\section{Cosets in the free-product join}

For the remainder of the argument let \(G=A*B\), where \(A,B\) are
nontrivial torsion-free groups. The factors are distinguished by their
inclusions even if they are abstractly isomorphic.
The free-product tree has vertices \(G/A\sqcup G/B\) and an edge \(g\)
joining \(gA\) to \(gB\), with \(G\) acting by left translation.
An element is called \emph{hyperbolic} if it fixes no vertex of this tree.

\begin{lemma}\label{lem:tree}
The conjugates of \(A,B\), together with the maximal cyclic subgroups
generated by hyperbolic elements of the free-product tree, satisfy:
\begin{enumerate}
\item two distinct subgroups in this collection intersect trivially;
\item the normalizer of every subgroup in the collection is itself;
\item the centralizer of a nonidentity element is contained in one
      subgroup in the collection.
\end{enumerate}
\end{lemma}

\begin{proof}
Reduced alternating words show that the free-product graph is
connected and has no circuits. Edge stabilizers are trivial; vertex
stabilizers are conjugates of the factors; the action preserves the
two vertex types and has no inversions.

A nonidentity element that fixes a vertex fixes no other vertex,
since otherwise it fixes an edge in the segment between them.
Its centralizer preserves its unique fixed vertex and lies in that
vertex stabilizer. Distinct vertex stabilizers therefore intersect
trivially. A normalizer of a vertex stabilizer preserves that vertex:
choose any nonidentity element of the stabilizer and use uniqueness
of its fixed vertex. This proves the normalizer assertion for factors.

A hyperbolic element \(g\) translates a unique axis, a bi-infinite geodesic.
To see this, choose a vertex \(v\) minimizing \(\ell=d(v,gv)>0\).
If \([v,gv]\) and \([gv,g^2v]\) backtrack at \(gv\), let \(w\) be the
neighbor of \(v\) toward \(gv\). Then \(gw\) is the neighbor of \(gv\)
toward \(v\). For \(\ell\geq2\) this gives \(d(w,gw)=\ell-2\), contrary
to minimality. For \(\ell=1\), the element exchanges the endpoints of
an edge, contrary to the absence of inversions. The translates therefore
concatenate without backtracking and give the axis, characterized by
minimum displacement, which also establishes its uniqueness.

The stabilizer of an axis acts faithfully on it, since an element in
the kernel fixes an edge. A reflection has square in this kernel and
would be an element of order two. Torsion freeness excludes it.
The axis stabilizer is thus a nonzero subgroup of the translations of
a simplicial line, so is infinite cyclic. It is the centralizer of
every nonidentity translation on that axis and the maximal cyclic
subgroup containing it.

Two such cyclic groups sharing a nonidentity element have the same
axis and equal its stabilizer. If \(hCh^{-1}\) meets \(C\) nontrivially,
their axes agree, so \(h\) belongs to the axis stabilizer \(C\).
This proves malnormality and the normalizer assertion for these
cyclic groups. They meet no vertex stabilizer nontrivially because
their nonidentity elements are translations. All assertions follow.
\end{proof}

Let \(\mathcal H\) contain \(A,B\) and one representative of each conjugacy
class of maximal hyperbolic cyclic subgroups. Call a subset
\[
 D=uHv^{-1}\subset G,\qquad H\in\mathcal H,\quad u,v\in G,
\]
a \emph{block}, and identify duplicate subsets. Let \(\mathcal B\) be the
set of blocks, and put
\[
 Z=\bigcup_{D\in\mathcal B}E(D)\subset E.
\]
For \(g\in G\), write \(F_g=E(\{g\})\). This is a contractible infinite
simplex, consisting of all simplices with the same group label \(g\)
at their different levels. Every \(F_g\) lies in \(Z\).

\begin{lemma}\label{lem:blocks}
Two distinct blocks intersect in at most one group element. Every cell
of \(E\) outside \(Z\) has trivial \(P\)-stabilizer. The setwise stabilizer
of the block \(H\) is \(H\times H\); it acts transitively on the labels of
that block, with stabilizer \(\Delta H\) at \(1\). The stabilizer of
the label \(1\) in all of \(G\) is \(\Delta G\).
\end{lemma}

\begin{proof}
A block is a left coset of a conjugate of \(H\). If two blocks contain
distinct \(g,h\), then \(g^{-1}h\) is a nonidentity element of both
associated subgroups. Lemma~\ref{lem:tree} identifies the subgroups,
and the common element \(g\) identifies their cosets.

For a simplex with nontrivial stabilizer, send one label to \(1\) by
the left-right action. A stabilizer element then has the form \((w,w)\),
where \(w\) centralizes all its labels. A nonidentity such \(w\) has
centralizer in a factor conjugate or an axis cyclic group, by
Lemma~\ref{lem:tree}. All the labels lie in that subgroup, so the
original simplex lies in one block. Constant-label simplices also
lie in \(Z\).

If \(uHv^{-1}=H\), then \(uv^{-1}=h\in H\). Substituting \(u=hv\)
shows that \(v\) normalizes \(H\), so \(u,v\in H\).
The converse and the asserted transitivity are immediate from the
action. Stabilizing the label \(1\) means \(u=v\), giving the two
diagonal subgroups as stated.
\end{proof}

\section{The mixed top classes survive}

We first record the relative comparison used to pass from Borel homology
to orbit homology.

\begin{lemma}\label{lem:relative}
For the pair \((E,Z)\) just constructed, the projection induces an
isomorphism
\[
 H_*(E_{hP},Z_{hP})\longrightarrow H_*(E/P,Z/P).
\]
\end{lemma}

\begin{proof}
Filter the pairs by \(Z\) together with the cell skeleta of \(E\).
By Lemma~\ref{lem:blocks}, each relative cell orbit is free.
For a \(k\)-cell its Borel relative piece is
\(EP\times(D^k,S^{k-1})\), and its orbit piece is
\((D^k,S^{k-1})\). Their projection is an isomorphism on relative
homology because \(EP\) is contractible.

For completeness, apply the same filtration to the mapping cone
of the cellular chain map induced by projection. Each associated
graded complex is a direct sum of mapping cones of the free-cell
projections above, hence is acyclic. A cycle of fixed total degree
has finite cellular support. Its component in the largest
filtration degree is a boundary in that graded complex.
Subtracting a lift of that boundary lowers the largest degree.
After finitely many repetitions the cycle is filled. The filtration
is bounded below, and the cell degrees and \(EP\)-chain degrees are
nonnegative, so no unbounded filtration enters a fixed total degree.
The mapping cone is acyclic, proving the homology isomorphism.
\end{proof}

Assume now that \(A,B\) admit classifying spaces with homology zero in
degrees greater than \(n\), where \(n\geq3\). In \(H_{2n}(P)\) define
\[
 V=\im H_{2n}(A\times A)+\im H_{2n}(B\times B),
\]
where the maps are induced by the two inclusions into \(G\times G\).

\begin{proposition}\label{prop:image}
The image of \(H_{2n}(Z_{hP})\) in \(H_{2n}(P)\) is contained in \(V\).
\end{proposition}

\begin{proof}
The block intersections give a \(P\)-equivariant pushout
\begin{equation}\label{eq:pushout}
\begin{tikzcd}[column sep=large,row sep=large]
 \displaystyle\bigsq_{D\in\mathcal B}\bigsq_{g\in D} F_g
       \arrow[r] \arrow[d] &
 \displaystyle\bigsq_{D\in\mathcal B}E(D) \arrow[d]\\
 \displaystyle\bigsq_{g\in G}F_g \arrow[r] & Z .
\end{tikzcd}
\end{equation}
Indeed, a simplex of \(Z\) with two distinct labels lies in exactly one
block. Distinct blocks are identified precisely along the \(F_g\)
they share. The top arrow is a CW cofibration, since in each block
the union of its constant-label subcomplexes is a subcomplex.
Thus this is a homotopy pushout.

One can take its explicit equivariant mapping-cylinder model before
forming Borel constructions. Products with \(EP\) and passage to
\(P\)-orbits commute with these identifications, and the subcomplex
inclusion remains a cofibration. The result is an ordinary homotopy
pushout. The blocks have one \(P\)-orbit for each \(H\in\mathcal H\),
with stabilizer \(H\times H\). Incidences within a representative
block form one orbit with stabilizer \(\Delta H\), and all labels
form one orbit with stabilizer \(\Delta G\). Since the joins in
the three input spaces are contractible, their Borel homotopy
types are
\begin{equation}\label{eq:borel-pushout}
\begin{tikzcd}[column sep=large,row sep=large]
 \displaystyle\bigsq_{H\in\mathcal H} B(\Delta H)
       \arrow[r] \arrow[d] &
 \displaystyle\bigsq_{H\in\mathcal H} B(H\times H) \arrow[d]\\
 B(\Delta G) \arrow[r] & Z_{hP}.
\end{tikzcd}
\end{equation}
Only these homotopy types are asserted in \eqref{eq:borel-pushout}.
Their maps to \(BP\) are the classifying maps of the corresponding
subgroup inclusions, up to conjugation.

Put \(d=2n\). For \(H=A,B\), the assumed homology bound gives
\(H_{d-1}(B\Delta H)=0\), because \(2n-1>n\).
For the remaining \(H\), the classifying space \(BH\) is a circle,
so the same vanishing holds. Also
\(H_d(B\Delta G)=0\): a classifying space for the free product is
\(BA\vee BB\), with no homology above \(n\).
Mayer--Vietoris for \eqref{eq:borel-pushout} therefore gives a surjection
\[
 \bigoplus_{H\in\mathcal H}H_d(H\times H)
       \longrightarrow H_d(Z_{hP}).
\]
The cyclic terms vanish in degree \(d\geq6\), since their product
classifying spaces are two-tori. The only possible contributions
come from \(A\times A\) and \(B\times B\), with exactly the inclusion
maps defining \(V\). Inner conjugation induces a homotopically
trivial change of a classifying-space map and does not enlarge
their images. This proves the proposition.
\end{proof}

\begin{proposition}\label{prop:survival}
If \(z\in H_{2n}(P)\) does not belong to \(V\), then
\(\tau_{2n}(z)\neq0\).
\end{proposition}

\begin{proof}
The exact homology sequence of \((E_{hP},Z_{hP})\), together with
Proposition~\ref{prop:image}, shows that \((\pi_*)^{-1}z\) has a
nonzero relative image. The commutative diagram
\[
\begin{tikzcd}[column sep=large,row sep=large]
 H_{2n}(E_{hP}) \arrow[r] \arrow[d,"\rho_*"] &
 H_{2n}(E_{hP},Z_{hP}) \arrow[d,"\cong"]\\
 H_{2n}(E/P) \arrow[r] & H_{2n}(E/P,Z/P)
\end{tikzcd}
\]
has an isomorphism on the right by Lemma~\ref{lem:relative}.
Thus \(\rho_*(\pi_*)^{-1}z\) has a nonzero relative image as well,
and is nonzero in absolute homology.
\end{proof}

\section{The connected-sum fundamental class}

\begin{proof}[Proof of Theorem~\ref{thm:main}]
Put \(X=M\#N\), \(A=\pi_1(M)\), \(B=\pi_1(N)\).
Closed topological manifolds have finite CW type, and the fundamental
group of an aspherical finite-dimensional CW complex is torsion free;
see~\cite[Theorem 5.3 and Lemma 3.1]{Luck}. Each factor is nontrivial:
a closed oriented contractible manifold would have both zero top
homology and a nonzero fundamental class. The free product \(G=A*B\)
is torsion free, since an element of finite order is conjugate into
a factor. In reduced-word terms, every other element is conjugate
to a cyclically reduced alternating word of length at least two,
whose positive powers have increasing length.

The wedge \(M\vee N\) is a classifying space for \(G\).
On based CW models its universal cover is a tree of the contractible
universal covers of the summands. Attaching points are cofibrations.
A finite subtree of contractible vertex spaces is contractible by
successive homotopy pushouts. A sphere map into the entire cover
has compact image in a finite CW subcomplex, so meets only finitely
many pieces and is null-homotopic in a finite subtree. Whitehead's
theorem makes the full cover contractible.

For \(n\geq3\), pinching the connecting sphere gives a classifying
map
\[
 c:X\longrightarrow M\vee N\simeq BG
\]
inducing the free-product fundamental-group isomorphism.
This is the usual connected-sum pinch map; compare the proof of
Lemma 3.2 in~\cite{Luck}. One may first collapse the neck and then
identify each punctured summand with its boundary collapsed with
the closed summand up to homotopy. The latter operation collapses
a contractible closed-ball cofibration.

Let \(a,b\in H_n(BG)\) be the separate images of \([M]\) and \([N]\).
Both projections of \(c\) onto a summand have degree \(+1\):
away from the removed ball each is the oriented identity near a
point with a single preimage, hence sends its local orientation
class to the positive generator. The two projections determine
the top homology of the wedge, so
\[
 c_*[X]=a+b.
\]
The rational K\"unneth theorem~\cite[Chapter 17, Section 2]{May} gives
\[
 H_{2n}(P)
   =\Q\{a\times a,\ a\times b,\ b\times a,\ b\times b\},
 \qquad V=\Q\{a\times a,\ b\times b\},
\]
since the homology of each summand vanishes above \(n\).
Naturality of the cross product yields
\[
 (c\times c)_*[X\times X]
   =a\times a+a\times b+b\times a+b\times b\notin V.
\]
The two mixed terms are independent ordered coordinates. They
remain distinct also when \(M=N\) or \(n\) is odd.

Proposition~\ref{prop:survival} now gives
\[
 \tau_{2n}\big((c\times c)_*[X\times X]\big)\neq0.
\]
If \(\TC(X)\leq2n\), Lemma~\ref{lem:planner} gives the opposite
conclusion in degree \(2n\). Hence \(\TC(X)\geq2n+1\).
Farber's dimension bound~\cite[Theorem 4]{Farber} gives the reverse
inequality, completing the proof.
\end{proof}

\paragraph{AI use.}
OpenAI Codex assisted with literature retrieval, mathematical development,
checking and manuscript preparation. Independent mathematical review is
requested.

\begin{thebibliography}{9}

\bibitem{Farber}
M. Farber, \emph{Topological complexity of motion planning},
Discrete \& Computational Geometry \textbf{29} (2003), 211--221.
\href{https://doi.org/10.1007/s00454-002-0760-9}{doi:10.1007/s00454-002-0760-9}.
Preprint \href{https://arxiv.org/abs/math/0111197v1}{arXiv:math/0111197v1}.

\bibitem{DSfree}
A. Dranishnikov and R. Sadykov,
\emph{The topological complexity of the free product},
\href{https://arxiv.org/abs/1710.04711v2}{arXiv:1710.04711v2} (2017).

\bibitem{Neofytidis}
C. Neofytidis, \emph{Topological complexity, asphericity and connected sums},
\href{https://arxiv.org/abs/2212.08962v2}{arXiv:2212.08962v2} (2023).

\bibitem{DST}
N. Daundkar, R. Santhanam and S. Thandar,
\emph{Higher topological complexity of Seifert fibered manifolds},
\href{https://arxiv.org/abs/2304.01274v3}{arXiv:2304.01274v3} (2026).

\bibitem{FGLO}
M. Farber, M. Grant, G. Lupton and J. Oprea,
\emph{Bredon cohomology and robot motion planning},
\href{https://arxiv.org/abs/1711.10132v2}{arXiv:1711.10132v2} (2018).

\bibitem{Luck}
W. L\"uck, \emph{Survey on aspherical manifolds},
\href{https://arxiv.org/abs/0902.2480v3}{arXiv:0902.2480v3} (2009).

\bibitem{May}
J. P. May, \emph{A Concise Course in Algebraic Topology},
University of Chicago Press, 1999.
\href{https://www.math.uchicago.edu/~may/CONCISE/ConciseRevised.pdf}{Author's revised text}.

\end{thebibliography}
\end{document}
